\subsection{Conorm of a continuous linear map}

Let E and F be normed spaces, and let \(u: E \rightarrow F\) be a continuous linear map, N the kernel of u and I its image. Let p denote the canonical surjection from E onto E/N and let i denote the canonical injection from I into F. Let \(\widetilde{u}\) be the bijective linear map from E/N onto I such that \(u = i \circ \widetilde{u} \circ p\) . The vector space E/N is equipped with the quotient norm, that is,

\[
\| y \| = \inf _ {x \in p^{-1}(y)} \| x \|\tag{2}
\]

for all \(y \in E/N\) . The map \(\widetilde{u}\) is continuous and \(\|u\| = \|\widetilde{u}\|\) , whence

\[
\| u\| = \sup_{\substack{y\in \mathrm{E / N}\\ y\neq 0}}\frac{\|\widetilde{u} (y)\|}{\|y\|},\tag{3}
\]

the supremum being taken in \(\overline{R}_{+}\) . The conorm of u is called the number

\[
((u)) = \inf_{\substack{y\in \mathrm{E} / \mathrm{N}\\ y\neq 0}}\frac{\| \widetilde{u}(y)\|}{\|y\|},\tag{4}
\]

the lower bound being taken in \(\overline{R}_{+}\) . We therefore have

\[
((u)) \| y \| \leqslant \| \widetilde {u} (y) \| \leqslant \| u \| \| y \|\tag{5}
\]

for every element \(y\) of E/N. Set \(v = i \circ \widetilde{u}\) . We have \(u = v \circ p\) and

\[
((u)) = \inf_{\substack{y\in \mathrm{E / N}\\ y\neq 0}}\frac{\|v(y)\|}{\|y\|} = \inf_{\substack{y\in \mathrm{E / N}\\ \|y\| = 1}}\| v(y)\| .\tag{6}
\]

When \(u\) is the zero map, the space E/N is reduced to 0 and one has \(((u)) = +\infty\) and \(\| u \| = 0\) . When \(u \neq 0\) , from (3) and (4) one deduces the inequalities

\[
0 \leqslant ((u)) \leqslant \| u \| <   + \infty .\tag{7}
\]

When u is injective, we have

\[
((u)) = \inf_{\substack{x\in \mathrm{E}\\ x\neq 0}}\frac{\|u(x)\|}{\|x\|},\tag{8}
\]

and, for all \(x \in \mathrm{E}\) , we have

\[
((u)) \| x \| \leqslant \| u (x) \| \leqslant \| u \| \| x \|.\tag{9}
\]

Denote by \(j\) the canonical injection of the normed space F into its completion \(\widehat{\mathbf{F}}\) . We have \(((u)) = ((j \circ u))\) .

Remark. --- By definition, \(((u)) > 0\) if and only if the bijective linear map \(\widetilde{u}\) is bicontinuous, that is, if and only if u is a strict morphism (TG, III, p. 16). We then have

\[
((u)) = \| \widetilde {u} ^ {- 1} \| ^ {- 1}.\tag{10}
\]

Lemma 2. --- Let c be a real number. If c \textless{} ((u)), then for every element \(z \in \operatorname{Im}(u)\) there exists an element x of E such that \(u(x) = z\) and \(c \|x\| \leqslant \|z\|\) . Conversely, if for every \(z \in \operatorname{Im}(u)\) there exists \(x \in E\) such that \(u(x) = z\) and \(c \|x\| \leqslant \|z\|\) , then \(c \leqslant ((u))\) .

This is a consequence of formulas (2) and (5) and of the definition of the conorm of \(u\) .

PROPOSITION 6. --- Let E and F be normed spaces and \(u \in \mathcal{L}(\mathrm{E};\mathrm{F})\) . Let B denote the set of elements of E with norm \(< 1\) . Set

\[
\mathrm{P} = u (\mathrm{E}) - u (\mathrm{B}), \quad \mathrm{Q} = u (\mathrm{E}) - (\overline {{u (\mathrm{B})}} \cap u (\mathrm{E})).
\]

The conorm of u is equal to the distance from 0 to P in F. If the normed space E is complete or if u is a strict morphism, the conorm of u is equal to the distance from 0 to Q in F.

Let N be the kernel of u; let \(p: E \to E/N\) be the canonical surjection, and let \(v: E/N \to F\) be the map induced by u on the quotient. The set \(p(B)\) is the set of elements of E/N of norm \textless{} 1. We have

\[
\mathrm{P} = u (\mathrm{E}) - u (\mathrm{B}) = v (\mathrm{E} / \mathrm{N}) - v (p (\mathrm{B})) = v ((\mathrm{E} / \mathrm{N}) - p (\mathrm{B}))
\]

since the map \(v\) is injective. In other words, P consists of the elements of F of the form \(v(y)\) with \(y \in \mathrm{E/N}\) and \(\|y\| \geqslant 1\) . Denote by \(d_{\mathrm{P}}\) the distance from 0 to P in F. We have

\[
d_{\mathrm{P}} = \inf_{\substack{y\in \mathrm{E} / \mathrm{N}\\ \| y\| \geqslant 1}}\| v(y)\| = \inf_{\substack{y\in \mathrm{E} / \mathrm{N}\\ \| y\| = 1}}\| v(y)\| = ((u))
\]

according to (6).

Suppose that \(u\) is a strict morphism. Let \(\varepsilon > 0\) . The set \(\varepsilon u(\mathrm{B})\) is a neighborhood of 0 in \(u(\mathrm{E})\) . The closure of \(u(\mathrm{B})\) in \(u(\mathrm{E})\) is equal to \(\overline{u(\mathrm{B})} \cap u(\mathrm{E})\) . It is contained in the set \(u(\mathrm{B}) + \varepsilon u(\mathrm{B})\) which is equal to \((1 + \varepsilon)u(\mathrm{B})\) since \(u(\mathrm{B})\) is convex. Consequently, we have \((1 + \varepsilon)\mathrm{P} \subset \mathrm{Q} \subset \mathrm{P}\) and the distance \(d_{\mathrm{Q}}\) from 0 to Q in F satisfies the inequalities \(d_{\mathrm{P}} \leqslant d_{\mathrm{Q}} \leqslant (1 + \varepsilon)d_{\mathrm{P}}\) . Since this holds for every \(\varepsilon > 0\) , we have \(d_{\mathrm{Q}} = d_{\mathrm{P}} = ((u))\) .

Suppose that \(u\) is not a strict morphism, but that the normed space E is complete. We then have \(((u)) = 0\) . The closure of \(u(\mathrm{B})\) in \(u(\mathrm{E})\) , which is equal to \(\overline{u(\mathrm{B})} \cap u(\mathrm{E})\) , is not a neighborhood of 0 in \(u(\mathrm{E})\) (EVT, I, p. 17, th. 1). There therefore exist points of Q arbitrarily close to 0, whence \(d_{\mathrm{Q}} = 0 = ((u))\) .

PROPOSITION 7. --- Let E be a Banach space and F a normed space. The mapping \(u \mapsto ((u))\) of \(\mathcal{L}(\mathrm{E};\mathrm{F})\) in \(\overline{\mathbf{R}}\) is upper semi-continuous.

Let \(u \in \mathcal{L}(\mathrm{E};\mathrm{F})\) . It suffices to prove that for every real number \(c > ((u))\) , the set of elements \(v \in \mathcal{L}(\mathrm{E};\mathrm{F})\) such that \(((v)) < c\) is a neighborhood of \(u\) . Let B denote the set of elements of E with norm \(< 1\) . By prop. 6, there exists \(y \in \mathrm{E}\) such that \(u(y) \notin \overline{u(\mathrm{B})}\) and \(\| u(y)\| < c\) . The distance \(d\) of \(u(y)\) to the closed set \(\overline{u(\mathrm{B})}\) is strictly positive. The set V of elements \(v\) of \(\mathcal{L}(\mathrm{E};\mathrm{F})\) satisfying the relations \(\| v(y)\| < c\) and \(\| u - v\|(1 + \| y\|) < d\) is a neighborhood of \(u\) in \(\mathcal{L}(\mathrm{E};\mathrm{F})\) . Let \(v \in \mathrm{V}\) . For every \(x \in \mathrm{B}\) , we have

\[
\begin{array}{c} d \leqslant \| u (y) - u (x) \| \leqslant \| v (y) - v (x) \| + \| u - v \| (\| y \| + \| x \|) \\ <   \| v (y) - v (x) \| + d. \end{array}
\]

Consequently \(v(y)\) does not belong to \(v(\mathrm{B})\) , and one has \(((v)) \leqslant \| v(y) \| < c\) by prop. 6.

PROPOSITION 8. --- Let E be a Banach space, F a normed space. For every \(u \in \mathcal{L}(\mathrm{E};\mathrm{F})\) , we have \(((u)) = ((^t u))\) .

Denote by \(j\) the canonical injection of the normed space F into its completion \(\widehat{\mathrm{F}}\) . We have \(((u)) = ((j \circ u))\) . Since the linear map \(^t j: (\widehat{\mathrm{F}})' \to \mathrm{F}'\) is bijective and isometric, one also has \(((^t u)) = ((^t (j \circ u)))\) . It therefore suffices to prove the proposition when the normed space F is complete, which we assume in the rest of the proof.

If \(u\) is zero, one has \(((u)) = ((^t u)) = +\infty\) . If \(u\) is not a strict morphism, then \(^t u\) is not one either (EVT, IV, p. 29, cor. 3), and one has \(((u)) = ((^t u)) = 0\) .

Suppose now that u is a nonzero strict morphism. Let N denote the kernel of u and I its image, and consider the canonical decomposition of u:

\[
\mathrm{E} \xrightarrow {p} \mathrm{E/N} \xrightarrow {v} \mathrm{I} \xrightarrow {i} \mathrm{F}.
\]

The kernel of \(^{t}u\) is the orthogonal I° of I in F' (EVT, IV, p. 27, prop. 2), and \(^{t}i\) induces by passage to the quotient an isometry \(\iota\) from F'/I° onto I' (EVT, IV, p. 9, prop. 10). Moreover, \(^{t}p\) defines an isometry from (E/N)' onto the orthogonal N° of N in E' (EVT, IV, p. 8, prop. 9). The canonical decomposition of \(^{t}u\) is therefore

\[
\mathrm{F} ^ {\prime} \longrightarrow \mathrm{F} ^ {\prime} / \mathrm{I} ^ {\circ} \stackrel {v ^ {\prime}} {\longrightarrow} \mathrm{N} ^ {\circ} \longrightarrow \mathrm{E} ^ {\prime},
\]

where \(v^{\prime} = \pi \circ{}^{t}v\circ \iota\). We then have

\[
\| (v ^ {\prime}) ^ {- 1} \| = \| (^ {t} v) ^ {- 1} \| = \| ^ {t} (v ^ {- 1}) \| = \| v ^ {- 1} \|
\]

(EVT, IV, p. 7, prop. 8), whence \(((u)) = ((^t u))\) according to formula (10).

